\documentclass[12pt,a4paper]{article}

% ===========================================
% PACKAGES ET CONFIGURATION
% ===========================================

\usepackage[utf8]{inputenc}
\usepackage[T1]{fontenc}
\usepackage[french]{babel}
\usepackage{geometry}
\usepackage{graphicx}
\usepackage{xcolor}
\usepackage{tikz}
\usetikzlibrary{shapes,arrows,positioning,fit,backgrounds}
\usepackage{listings}
\usepackage{fancyhdr}
\usepackage{tocloft}
\usepackage{hyperref}
\usepackage{enumitem}
\usepackage{titlesec}
\usepackage{caption}
\usepackage{subcaption}
\usepackage{amsmath}
\usepackage{amssymb}
\usepackage{booktabs}
\usepackage{longtable}
\usepackage{array}
\usepackage{tabularx}
\usepackage{float}

% Configuration identique à la partie 1
\geometry{
    left=2.5cm,
    right=2.5cm,
    top=3cm,
    bottom=3cm,
    headheight=15pt
}

% Configuration des couleurs
\definecolor{darkgray}{RGB}{44,44,44}
\definecolor{lightblue}{RGB}{100,181,246}
\definecolor{lightpurple}{RGB}{186,104,200}
\definecolor{lightgreen}{RGB}{129,199,132}
\definecolor{lightorange}{RGB}{255,183,77}
\definecolor{lightpink}{RGB}{240,98,146}
\definecolor{lightcyan}{RGB}{66,165,245}
\definecolor{orange}{RGB}{255,167,38}

% Configuration des listings
\lstdefinestyle{bash}{
    language=bash,
    basicstyle=\ttfamily\footnotesize,
    backgroundcolor=\color{gray!10},
    frame=single,
    breaklines=true,
    showstringspaces=false,
    commentstyle=\color{green!60!black},
    keywordstyle=\color{blue!80!black},
    stringstyle=\color{red!80!black},
    numbers=left,
    numberstyle=\tiny\color{gray},
    stepnumber=1,
    numbersep=8pt
}

\lstdefinestyle{python}{
    language=Python,
    basicstyle=\ttfamily\footnotesize,
    backgroundcolor=\color{gray!10},
    frame=single,
    breaklines=true,
    showstringspaces=false,
    commentstyle=\color{green!60!black},
    keywordstyle=\color{blue!80!black},
    stringstyle=\color{red!80!black},
    numbers=left,
    numberstyle=\tiny\color{gray},
    stepnumber=1,
    numbersep=8pt
}

\lstdefinestyle{javascript}{
    language=JavaScript,
    basicstyle=\ttfamily\footnotesize,
    backgroundcolor=\color{gray!10},
    frame=single,
    breaklines=true,
    showstringspaces=false,
    commentstyle=\color{green!60!black},
    keywordstyle=\color{blue!80!black},
    stringstyle=\color{red!80!black},
    numbers=left,
    numberstyle=\tiny\color{gray},
    stepnumber=1,
    numbersep=8pt
}

% Configuration des titres
\titleformat{\section}
  {\Large\bfseries\color{blue!80!black}}
  {\thesection}
  {1em}
  {}
  [\color{blue!80!black}\titlerule]

\titleformat{\subsection}
  {\large\bfseries\color{blue!70!black}}
  {\thesubsection}
  {1em}
  {}

\titleformat{\subsubsection}
  {\normalsize\bfseries\color{blue!60!black}}
  {\thesubsubsection}
  {1em}
  {}

% Configuration des headers/footers
\pagestyle{fancy}
\fancyhf{}
\fancyhead[L]{\leftmark}
\fancyhead[R]{Documentation Technique Doc@uthANTIC - Partie 2}
\fancyfoot[C]{\thepage}
\renewcommand{\headrulewidth}{0.4pt}
\renewcommand{\footrulewidth}{0pt}

% Configuration des liens
\hypersetup{
    colorlinks=true,
    linkcolor=blue!80!black,
    filecolor=magenta,
    urlcolor=cyan!80!black,
    citecolor=green!80!black,
}

\begin{document}

\title{Documentation Technique Doc@uthANTIC\\Implémentation Nginx Reverse Proxy\\{\large Partie 2 : Résolution des problèmes et architecture finale}}
\author{Équipe de développement Doc@uthANTIC}
\date{\today}
\maketitle

\tableofcontents
\newpage

% ===========================================
% SECTION 7: CORRECTION PERMISSIONS DJANGO
% ===========================================

\section{Correction des permissions Django}

\subsection{Problème de permissions identifié}

Après la résolution des problèmes CORS et CSRF, un nouveau problème est apparu : les signataires (\textit{signers}) ne pouvaient pas accéder aux documents de leur organisation. L'erreur suivante était retournée :

\begin{lstlisting}[style=bash, caption={Erreur de permissions pour les signers}]
GET /api/documents/qr-positions/pending_for_signer/?organization_id=2 HTTP/1.0" 200 -
# Retournait une liste vide [] au lieu des documents attendus
\end{lstlisting}

\subsection{Analyse du problème}

\subsubsection{Code problématique dans DocumentQRPositionViewSet}

\begin{lstlisting}[style=python, caption={Code original avec problème de permissions}]
# backend/django-project/documents/views.py - AVANT

class DocumentQRPositionViewSet(viewsets.ModelViewSet):
    queryset = DocumentQRPosition.objects.all()
    serializer_class = DocumentQRPositionSerializer
    permission_classes = [IsAuthenticated]
    filter_backends = [DjangoFilterBackend, OrderingFilter]
    filterset_fields = ['status', 'signer']
    ordering_fields = ['created_at', 'updated_at']
    ordering = ['-created_at']

    def get_queryset(self):
        user = self.request.user
        organization_id = self.request.query_params.get('organization_id')
        
        # PROBLÈME : Logique restrictive pour les signers
        if user.role == 'signer':
            # ❌ Filtrait uniquement par signer = user
            return DocumentQRPosition.objects.filter(
                signer=user,
                organization_id=organization_id
            )
        else:
            # Pour admin et collaborateur
            return DocumentQRPosition.objects.filter(
                organization_id=organization_id
            )
\end{lstlisting}

\subsubsection{Origine du problème}
\begin{itemize}[leftmargin=*]
    \item \textbf{Filtrage trop restrictif :} Un signer ne voyait que ses propres documents
    \item \textbf{Logic business incorrecte :} Les signers doivent voir tous les documents de l'organisation
    \item \textbf{Workflow de signature bloqué :} Impossible d'accéder aux documents à signer
\end{itemize}

\subsection{Solution implémentée}

\subsubsection{Code corrigé avec permissions appropriées}

\begin{lstlisting}[style=python, caption={Code corrigé pour les permissions des signers}]
# backend/django-project/documents/views.py - APRÈS

class DocumentQRPositionViewSet(viewsets.ModelViewSet):
    queryset = DocumentQRPosition.objects.all()
    serializer_class = DocumentQRPositionSerializer
    permission_classes = [IsAuthenticated]
    filter_backends = [DjangoFilterBackend, OrderingFilter]
    filterset_fields = ['status', 'signer']
    ordering_fields = ['created_at', 'updated_at']
    ordering = ['-created_at']

    def get_queryset(self):
        user = self.request.user
        organization_id = self.request.query_params.get('organization_id')
        
        # ✅ CORRECTION : Tous les rôles peuvent voir tous les documents de l'organisation
        if organization_id:
            return DocumentQRPosition.objects.filter(
                organization_id=organization_id
            )
        
        # Fallback si pas d'organization_id
        if user.role == 'admin':
            return DocumentQRPosition.objects.all()
        else:
            # Pour collaborateur et signer : documents de leur organisation
            return DocumentQRPosition.objects.filter(
                organization=user.organization
            )

    @action(detail=False, methods=['get'])
    def pending_for_signer(self, request):
        """
        ✅ NOUVEAU : Endpoint spécifique pour les documents en attente de signature
        """
        user = request.user
        organization_id = request.query_params.get('organization_id')
        
        if organization_id:
            # Filtrer les documents en attente dans l'organisation
            queryset = DocumentQRPosition.objects.filter(
                organization_id=organization_id,
                status='pending'  # Ou le statut approprié
            )
        else:
            queryset = DocumentQRPosition.objects.filter(
                organization=user.organization,
                status='pending'
            )
        
        serializer = self.get_serializer(queryset, many=True)
        return Response(serializer.data)
\end{lstlisting}

\subsection{Tests de validation des permissions}

\begin{lstlisting}[style=bash, caption={Tests des permissions corrigées}]
# Test : Accès aux documents par un signer
curl -k -H "Authorization: Bearer $JWT_TOKEN" \
     -H "Content-Type: application/json" \
     https://192.168.4.131/api/documents/qr-positions/?organization_id=2

# Réponse attendue : Liste de tous les documents de l'organisation 2
# [
#   {
#     "id": "3166e3fc-fade-4b8c-ab73-f7a9926c99d0",
#     "document_name": "nzaangwai_084151.pdf",
#     "status": "pending",
#     "organization": 2,
#     ...
#   }
# ]

# Test : Documents en attente de signature
curl -k -H "Authorization: Bearer $JWT_TOKEN" \
     https://192.168.4.131/api/documents/qr-positions/pending_for_signer/?organization_id=2

# Réponse attendue : Documents avec status="pending"
\end{lstlisting}

% ===========================================
% SECTION 8: PROBLÈMES DE RÉÉCRITURE D'URL
% ===========================================

\section{Problèmes de réécriture d'URL}

\subsection{Problème d'URL rewriting identifié}

Après avoir résolu les problèmes de permissions, un nouveau problème est apparu avec les URLs du microservice de signature. Les requêtes vers \texttt{/sign/} retournaient des erreurs 404.

\subsection{Analyse du problème de rewrite}

\subsubsection{Configuration Nginx problématique}

\begin{lstlisting}[style=bash, caption={Configuration Nginx causant le problème}]
# Configuration problématique dans /etc/nginx/sites-available/docuthanatic
location /sign/ {
    rewrite ^/sign/(.*)$ /$1 break;  # ❌ PROBLÈME ICI
    proxy_pass https://127.0.0.1:8003;
    # ... autres headers ...
}
\end{lstlisting}

\subsubsection{Origine du problème}
\begin{itemize}[leftmargin=*]
    \item \textbf{URL d'origine :} \texttt{https://192.168.4.131/sign/}
    \item \textbf{Après rewrite :} \texttt{/(chaîne vide)} → 404 sur le microservice
    \item \textbf{URL attendue par le microservice :} \texttt{/sign}
    \item \textbf{Solution :} Frontend doit envoyer \texttt{/sign/sign} pour obtenir \texttt{/sign}
\end{itemize}

\subsection{Solution frontend implémentée}

\subsubsection{Modifications des URLs dans les services Vue.js}

\begin{lstlisting}[style=javascript, caption={Correction des URLs dans le frontend}]
// frontend/src/services/DocumentService.js - AVANT
signDocument(documentId, signatureData) {
    const url = `/sign/process_signature/${documentId}/`;  // ❌ Devenait /process_signature/
    return apiClient.post(url, signatureData);
}

// frontend/src/services/DocumentService.js - APRÈS  
signDocument(documentId, signatureData) {
    const url = `/sign/sign/process_signature/${documentId}/`;  // ✅ Devient /sign/process_signature/
    return apiClient.post(url, signatureData);
}

// frontend/src/services/CryptoService.js - AVANT
generateCertificate(certificateData) {
    const url = `/cert/generate_certificate/`;  // ❌ Devenait /generate_certificate/
    return apiClient.post(url, certificateData);
}

// frontend/src/services/CryptoService.js - APRÈS
generateCertificate(certificateData) {
    const url = `/cert/cert/generate_certificate/`;  // ✅ Devient /cert/generate_certificate/
    return apiClient.post(url, certificateData);
}
\end{lstlisting}

\subsection{Table de correspondance URL finale}

\begin{table}[H]
\centering
\caption{Correspondance URLs frontend → microservices après correction}
\label{tab:url_mapping_final}
\begin{tabularx}{\textwidth}{|X|X|X|}
\hline
\textbf{URL Frontend} & \textbf{Après Nginx rewrite} & \textbf{Microservice reçoit} \\
\hline
\texttt{/sign/sign/process\_signature/} & \texttt{/sign/process\_signature/} & ✅ Correct \\
\hline
\texttt{/cert/cert/generate\_certificate/} & \texttt{/cert/generate\_certificate/} & ✅ Correct \\
\hline
\texttt{/gateway/gateway/health/} & \texttt{/gateway/health/} & ✅ Correct \\
\hline
\texttt{/api/users/auth/} & \texttt{/api/users/auth/} & ✅ Pas de rewrite \\
\hline
\end{tabularx}
\end{table}

\subsection{Tests de validation des URLs}

\begin{lstlisting}[style=bash, caption={Tests des URLs corrigées}]
# Test microservice signature
curl -k -X POST \
     -H "Content-Type: application/json" \
     -H "Authorization: Bearer $JWT_TOKEN" \
     -d '{"signature_data": "..."}' \
     https://192.168.4.131/sign/sign/process_signature/12345/

# Réponse attendue : 200 OK avec données de signature

# Test microservice certificat  
curl -k -X POST \
     -H "Content-Type: application/json" \
     -H "Authorization: Bearer $JWT_TOKEN" \
     -d '{"certificate_data": "..."}' \
     https://192.168.4.131/cert/cert/generate_certificate/

# Réponse attendue : 200 OK avec certificat généré
\end{lstlisting}

% ===========================================
% SECTION 9: ARCHITECTURE FINALE
% ===========================================

\section{Architecture finale et flux de données}

\subsection{Diagramme de l'architecture finale}

\begin{figure}[H]
    \centering
    \begin{tikzpicture}[
        node distance=1.5cm and 2cm,
        nginx/.style={rectangle, draw, fill=black, text=white, minimum width=3cm, minimum height=1.2cm, font=\small\bfseries},
        box/.style={rectangle, draw, fill=darkgray, text=white, minimum width=2.8cm, minimum height=1cm, font=\small\bfseries},
        client/.style={rectangle, draw, fill=blue!20, minimum width=2.5cm, minimum height=0.8cm, font=\small},
        db/.style={cylinder, draw, fill=lightcyan, text=black, minimum width=2cm, minimum height=1cm, font=\small\bfseries},
        arrow/.style={->, thick}
    ]
        % Client
        \node[client] (client) {Client Web Browser};
        
        % Nginx (point central)
        \node[nginx, below=1.5cm of client] (nginx) {Nginx Reverse Proxy\\192.168.4.131:443\\(Point d'entrée unique)};
        
        % Services backend (localhost)
        \node[box, below=2.5cm of nginx, xshift=-6cm] (frontend) {Vue.js Frontend\\127.0.0.1:8080};
        \node[box, right=1cm of frontend] (django) {Django Backend\\127.0.0.1:8000};
        \node[box, right=1cm of django] (gateway) {API Gateway\\127.0.0.1:8001};
        \node[box, below=1.5cm of frontend] (certms) {Certificat MS\\127.0.0.1:8002};
        \node[box, below=1.5cm of django] (signms) {Signature MS\\127.0.0.1:8003};
        
        % Base de données
        \node[db, below=1.5cm of gateway] (postgres) {PostgreSQL\\Database};
        
        % Connexions client → Nginx
        \draw[arrow, green!80!black, very thick] (client) -- node[right, font=\footnotesize] {HTTPS:443} (nginx);
        
        % Connexions Nginx → Services
        \draw[arrow, lightblue, thick] (nginx) to[out=-135,in=90] node[left, font=\tiny] {/:8080} (frontend);
        \draw[arrow, lightpurple, thick] (nginx) to[out=-90,in=90] node[above, font=\tiny] {/api/:8000} (django);
        \draw[arrow, lightgreen, thick] (nginx) to[out=-45,in=90] node[right, font=\tiny] {/gateway/:8001} (gateway);
        \draw[arrow, lightorange, thick] (nginx) to[out=-115,in=45] node[left, font=\tiny] {/cert/:8002} (certms);
        \draw[arrow, lightpink, thick] (nginx) to[out=-65,in=135] node[right, font=\tiny] {/sign/:8003} (signms);
        
        % Connexions inter-services
        \draw[arrow, dashed] (frontend) -- node[above, font=\tiny] {API calls} (django);
        \draw[arrow, dashed] (frontend) to[out=45,in=135] node[above, font=\tiny] {Gateway} (gateway);
        \draw[arrow, dashed] (gateway) -- node[left, font=\tiny] {Cert MS} (certms);
        \draw[arrow, dashed] (gateway) -- node[right, font=\tiny] {Sign MS} (signms);
        
        % Connexions base de données
        \draw[arrow, dashed] (django) to[out=-45,in=90] node[right, font=\tiny] {ORM} (postgres);
        \draw[arrow, dashed] (certms) to[out=0,in=180] node[below, font=\tiny] {Store} (postgres);
        \draw[arrow, dashed] (signms) to[out=180,in=0] node[below, font=\tiny] {Store} (postgres);
        
        % Titre
        \node[above=0.5cm of client, font=\Large\bfseries] {Architecture finale Doc@uthANTIC};
        
        % Légende
        \node[below=4cm of postgres, font=\small, text=green!80!black] {✅ Architecture sécurisée avec reverse proxy Nginx};
    \end{tikzpicture}
    \caption{Architecture finale Doc@uthANTIC avec Nginx reverse proxy}
    \label{fig:architecture_finale}
\end{figure}

\subsection{Flux de données complet}

\subsubsection{Diagramme de séquence : Signature d'un document}

\begin{figure}[H]
    \centering
    \begin{tikzpicture}[
        node distance=0.5cm and 1.5cm,
        participant/.style={rectangle, draw, fill=blue!10, minimum width=2cm, minimum height=0.8cm, font=\footnotesize},
        arrow/.style={->, thick}
    ]
        % Participants
        \node[participant] (user) {Utilisateur};
        \node[participant, right=2cm of user] (nginx) {Nginx};
        \node[participant, right=2cm of nginx] (frontend) {Frontend};
        \node[participant, right=2cm of frontend] (django) {Django};
        \node[participant, below=4cm of nginx] (gateway) {Gateway};
        \node[participant, below=4cm of frontend] (signms) {Sign MS};
        \node[participant, below=4cm of django] (db) {PostgreSQL};
        
        % Lignes de vie
        \draw[dashed] (user) -- ++(0,-8);
        \draw[dashed] (nginx) -- ++(0,-8);
        \draw[dashed] (frontend) -- ++(0,-8);
        \draw[dashed] (django) -- ++(0,-8);
        \draw[dashed] (gateway) -- ++(0,-8);
        \draw[dashed] (signms) -- ++(0,-8);
        \draw[dashed] (db) -- ++(0,-8);
        
        % Messages
        \draw[arrow] (user) -- ++(2,0) -- ++(0,-0.5) node[right, font=\tiny] {1. HTTPS Request};
        \draw[arrow] (nginx) -- ++(2,0) -- ++(0,-0.5) node[right, font=\tiny] {2. Route /};
        \draw[arrow] (frontend) -- ++(2,0) -- ++(0,-0.5) node[right, font=\tiny] {3. API /api/};
        \draw[arrow] (django) -- ++(2,0) -- ++(0,-1) node[right, font=\tiny] {4. Auth \& Data};
        \draw[arrow] (django) -- ++(0,-1.5) -- ++(-2,0) node[left, font=\tiny] {5. Call /sign/sign/};
        \draw[arrow] (gateway) -- ++(2,0) -- ++(0,-0.5) node[right, font=\tiny] {6. Process};
        \draw[arrow] (signms) -- ++(2,0) -- ++(0,-0.5) node[right, font=\tiny] {7. Store};
        
        % Retours
        \draw[arrow, dashed] (db) -- ++(-2,0) -- ++(0,-0.5) node[left, font=\tiny] {Response};
        \draw[arrow, dashed] (signms) -- ++(-2,0) -- ++(0,-0.5) node[left, font=\tiny] {Success};
        \draw[arrow, dashed] (gateway) -- ++(-2,0) -- ++(0,-0.5) node[left, font=\tiny] {Result};
        \draw[arrow, dashed] (django) -- ++(-2,0) -- ++(0,-0.5) node[left, font=\tiny] {JSON};
        \draw[arrow, dashed] (frontend) -- ++(-2,0) -- ++(0,-0.5) node[left, font=\tiny] {UI Update};
        \draw[arrow, dashed] (nginx) -- ++(-2,0) -- ++(0,-0.5) node[left, font=\tiny] {Final Response};
    \end{tikzpicture}
    \caption{Flux de signature d'un document via Nginx}
    \label{fig:flux_signature}
\end{figure}

\subsection{Points clés de l'architecture finale}

\subsubsection{✅ Sécurité}
\begin{itemize}[leftmargin=*]
    \item \textbf{Point d'entrée unique :} Nginx sur ports 80/443
    \item \textbf{Services en localhost :} Inaccessibles depuis l'extérieur
    \item \textbf{SSL centralisé :} Gestion des certificats sur Nginx uniquement
    \item \textbf{Headers de sécurité :} X-Frame-Options, X-XSS-Protection, etc.
\end{itemize}

\subsubsection{✅ Performance}
\begin{itemize}[leftmargin=*]
    \item \textbf{Load balancing possible :} Via Nginx upstream
    \item \textbf{Cache statique :} Nginx peut servir les assets
    \item \textbf{Compression :} Gzip activée sur Nginx
    \item \textbf{Keep-alive :} Connexions persistantes
\end{itemize}

\subsubsection{✅ Maintenance}
\begin{itemize}[leftmargin=*]
    \item \textbf{URLs simplifiées :} Plus de ports dans les URLs
    \item \textbf{Configuration centralisée :} Routage dans Nginx
    \item \textbf{Logs centralisés :} Accès et erreurs sur Nginx
    \item \textbf{Déploiement simplifié :} Un seul point d'entrée
\end{itemize}

% ===========================================
% SECTION 10: TESTS ET VALIDATION
% ===========================================

\section{Tests et validation}

\subsection{Suite de tests de connectivité}

\begin{lstlisting}[style=bash, caption={Tests de connectivité complets}]
#!/bin/bash
# Script de test complet de l'architecture Nginx

echo "=== Tests de connectivité Doc@uthANTIC avec Nginx ==="

# Test 1: Accès frontend
echo "1. Test frontend Vue.js..."
curl -k -s -o /dev/null -w "%{http_code}" https://192.168.4.131/
# Attendu: 200

# Test 2: API Django
echo "2. Test API Django..."
curl -k -s -o /dev/null -w "%{http_code}" https://192.168.4.131/api/
# Attendu: 200 ou 404 (pas d'endpoint racine)

# Test 3: Django Admin
echo "3. Test Django Admin..."
curl -k -s -o /dev/null -w "%{http_code}" https://192.168.4.131/admin/
# Attendu: 200 (redirection vers login)

# Test 4: API Gateway
echo "4. Test API Gateway..."
curl -k -s -o /dev/null -w "%{http_code}" https://192.168.4.131/gateway/
# Attendu: 200

# Test 5: Microservice Certificat
echo "5. Test Microservice Certificat..."
curl -k -s -o /dev/null -w "%{http_code}" https://192.168.4.131/cert/
# Attendu: 200

# Test 6: Microservice Signature
echo "6. Test Microservice Signature..."
curl -k -s -o /dev/null -w "%{http_code}" https://192.168.4.131/sign/
# Attendu: 200

echo "=== Tests terminés ==="
\end{lstlisting}

\subsection{Tests de sécurité}

\begin{lstlisting}[style=bash, caption={Tests de sécurité}]
# Test 1: Vérification que les services ne sont pas accessibles directement
echo "Test sécurité: Accès direct aux services (doit échouer)"

# Ces commandes doivent échouer (timeout ou connection refused)
curl -k --connect-timeout 5 https://192.168.4.131:8000/ || echo "✅ Django inaccessible directement"
curl -k --connect-timeout 5 https://192.168.4.131:8001/ || echo "✅ Gateway inaccessible directement"
curl -k --connect-timeout 5 https://192.168.4.131:8002/ || echo "✅ Cert MS inaccessible directement"
curl -k --connect-timeout 5 https://192.168.4.131:8003/ || echo "✅ Sign MS inaccessible directement"

# Test 2: Vérification des headers de sécurité
echo "Test headers de sécurité..."
curl -k -I https://192.168.4.131/ | grep -E "(X-Frame-Options|X-XSS-Protection|X-Content-Type-Options)"
# Attendu: Présence des headers de sécurité
\end{lstlisting}

\subsection{Tests fonctionnels}

\begin{lstlisting}[style=bash, caption={Tests fonctionnels complets}]
# Test d'authentification complète
echo "=== Test d'authentification ==="

# 1. Authentification
AUTH_RESPONSE=$(curl -k -s -X POST \
  -H "Content-Type: application/json" \
  -d '{
    "email": "test@example.com",
    "password": "testpass",
    "organization_id": 1
  }' \
  https://192.168.4.131/api/users/auth/with-organization/)

echo "Auth response: $AUTH_RESPONSE"

# 2. Extraction du token JWT
ACCESS_TOKEN=$(echo $AUTH_RESPONSE | grep -o '"access":"[^"]*"' | cut -d'"' -f4)
echo "Token: ${ACCESS_TOKEN:0:50}..."

# 3. Test d'accès aux documents avec token
echo "=== Test accès documents avec JWT ==="
curl -k -H "Authorization: Bearer $ACCESS_TOKEN" \
     https://192.168.4.131/api/documents/qr-positions/?organization_id=1

# 4. Test microservice signature avec token
echo "=== Test microservice signature ==="
curl -k -X POST \
     -H "Authorization: Bearer $ACCESS_TOKEN" \
     -H "Content-Type: application/json" \
     -d '{"test": "data"}' \
     https://192.168.4.131/sign/sign/health/
\end{lstlisting}

% ===========================================
% SECTION 11: RECOMMANDATIONS ET CONCLUSION
% ===========================================

\section{Recommandations et optimisations}

\subsection{Sécurité}

\subsubsection{Firewall et fail2ban}
\begin{lstlisting}[style=bash, caption={Configuration recommandée du firewall}]
# Configuration UFW (Uncomplicated Firewall)
sudo ufw enable
sudo ufw default deny incoming
sudo ufw default allow outgoing

# Autoriser seulement SSH et HTTPS
sudo ufw allow ssh
sudo ufw allow 443/tcp
sudo ufw allow 80/tcp

# Bloquer l'accès direct aux ports des services
sudo ufw deny 8000:8003/tcp

# Installation et configuration fail2ban
sudo apt install fail2ban -y
sudo systemctl enable fail2ban
\end{lstlisting}

\subsubsection{Rate limiting Nginx}
\begin{lstlisting}[style=bash, caption={Configuration rate limiting}]
# Ajouter dans /etc/nginx/sites-available/docuthanatic

# Zone de rate limiting
limit_req_zone $binary_remote_addr zone=api:10m rate=10r/s;
limit_req_zone $binary_remote_addr zone=login:10m rate=1r/s;

server {
    # Dans location /api/
    location /api/ {
        limit_req zone=api burst=20 nodelay;
        # ... reste de la configuration
    }
    
    # Rate limiting spécial pour login
    location /api/users/auth/ {
        limit_req zone=login burst=5 nodelay;
        # ... reste de la configuration
    }
}
\end{lstlisting}

\subsection{Performance}

\subsubsection{Cache et compression}
\begin{lstlisting}[style=bash, caption={Optimisations performance Nginx}]
# Configuration dans /etc/nginx/sites-available/docuthanatic

server {
    # Compression
    gzip on;
    gzip_vary on;
    gzip_min_length 1024;
    gzip_types
        text/plain
        text/css
        text/xml
        text/javascript
        application/javascript
        application/json
        application/xml+rss;
    
    # Cache pour assets statiques
    location ~* \.(js|css|png|jpg|jpeg|gif|ico|svg)$ {
        expires 1y;
        add_header Cache-Control "public, immutable";
        proxy_pass https://127.0.0.1:8000;
    }
    
    # Cache pour media files
    location /media/ {
        expires 30d;
        add_header Cache-Control "public";
        proxy_pass https://127.0.0.1:8000;
    }
}
\end{lstlisting}

\subsection{Monitoring}

\subsubsection{Logs et monitoring}
\begin{lstlisting}[style=bash, caption={Configuration monitoring}]
# Rotation des logs Nginx
sudo nano /etc/logrotate.d/nginx

/var/log/nginx/*.log {
    daily
    missingok
    rotate 52
    compress
    delaycompress
    notifempty
    create 644 nginx nginx
    sharedscripts
    prerotate
        if [ -d /etc/logrotate.d/httpd-prerotate ]; then \
            run-parts /etc/logrotate.d/httpd-prerotate; \
        fi
    endscript
    postrotate
        systemctl reload nginx
    endscript
}

# Monitoring en temps réel
sudo tail -f /var/log/nginx/docuthanatic_access.log
sudo tail -f /var/log/nginx/docuthanatic_error.log
\end{lstlisting}

\section{Conclusion}

\subsection{Résumé des accomplissements}

L'implémentation du reverse proxy Nginx pour Doc@uthANTIC a été un succès complet. Nous avons transformé une architecture exposant 5 services sur des ports distincts en une architecture sécurisée avec un point d'entrée unique.

\subsubsection{✅ Problèmes résolus}
\begin{enumerate}[leftmargin=*]
    \item \textbf{CORS :} Configuration Django mise à jour pour accepter les requêtes via Nginx
    \item \textbf{CSRF :} Transmission des tokens via headers Nginx et exemption JWT
    \item \textbf{Permissions :} Correction de la logique pour permettre aux signers d'accéder aux documents
    \item \textbf{URL Rewriting :} Adaptation des URLs frontend pour la réécriture Nginx
\end{enumerate}

\subsubsection{✅ Bénéfices obtenus}
\begin{enumerate}[leftmargin=*]
    \item \textbf{Sécurité :} Services en localhost, SSL centralisé
    \item \textbf{Simplicité :} URLs sans ports, configuration centralisée
    \item \textbf{Performance :} Possibilité de cache, compression, load balancing
    \item \textbf{Monitoring :} Logs centralisés sur Nginx
\end{enumerate}

\subsubsection{🎯 Recommandations finales}
\begin{enumerate}[leftmargin=*]
    \item \textbf{Monitoring :} Mettre en place Prometheus + Grafana
    \item \textbf{Load Balancing :} Configurer upstream pour haute disponibilité
    \item \textbf{Sécurité :} Rate limiting et fail2ban
    \item \textbf{Sauvegardes :} Automatisation des backups de configuration
\end{enumerate}

\textbf{🎉 Le système Doc@uthANTIC avec Nginx reverse proxy est maintenant en production et opérationnel !}

\subsection{Références}

\begin{itemize}[leftmargin=*]
    \item \href{https://nginx.org/en/docs/}{Nginx Documentation}
    \item \href{https://certisign.readthedocs.io/}{Doc@uthANTIC Documentation}
    \item \href{https://fastapi.tiangolo.com/}{Python FastAPI Documentation}
    \item \href{https://docs.djangoproject.com/}{Django Documentation}
    \item \href{https://vuejs.org/guide/}{Vue.js Documentation}
\end{itemize}

\end{document} 